\chapter*{Реферат}
\addcontentsline{toc}{chapter}{Реферат}

Выпускная квалификационная работа бакалавра содержит \pageref{LastPage} страниц, \ref{fig_count} рисунков, \ref{tab_count} таблиц, \ref{ref_count} источников.

\textbf{Ключевые слова:} уравнение Вебстера, нейросетевые операторы, оператор Фурье (FNO), DeepONet, Mamba, сверточная сеть UNet, прямая и обратная задачи волноводной акустики.

\textbf{Объект исследования:} одномерные акустические волноводы переменного сечения и процессы распространения звуковых волн в них.

\textbf{Цель работы:} исследование и сравнительный анализ нейросетевых архитектур для решения прямой и обратной задач в одномерной акустике волноводов переменного сечения: быстрого предсказания спектрального отклика по геометрическому профилю и восстановления геометрического профиля по спектральному отклику.

В ходе работы был реализован набор численных решателей прямой задачи (TMM, ARMA, TLM), использованный для генерации данных, верификации результатов и сравнения с нейросетевыми моделями. Реализованы и обучены нейросетевые аппроксиматоры прямой задачи, включая MLP, FNO, DeepONet, Mamba Operator и ряд модификаций DeepONet. Проведено исследование влияния различных архитектурных решений в DeepONet, включая механизмы Mamba, SIREN, FiLM, Attention Fusion, а также динамические и деформируемые свертки. Для обратной задачи восстановления геометрии по спектральному отклику сопоставлены оптимизационные и нейросетевые подходы, включая градиентную оптимизацию, метод Метрополиса и специализированную сверточную модель UNet, обученную на двумерных пространственно-спектральных картах входа $[H, f, x]$.
